\section{Experiments}
\label{sec:experiments}

We evaluate two claims separately. The budget experiments ask whether an information-weighted policy family can meet an expected sensing-usage target, transfer that calibration across reward scales, and re-adapt after a scale change. The EFE experiments ask where the canonical $w{=}1$ policy lies relative to reward-only planning, tuned information gain, and established POMDP solvers. Meeting a usage target and maximizing reward under a usage cap are different objectives, so the first battery measures both usage error and reward relative to a constrained reference.

The six main environments form a progression in sensing decisions. Tiger has one sensor and asks whether to listen again. Diagnosis and Bandit offer several tests and ask which observation to buy. Tileworld makes those tests spatial partitions of a grid. RockSample and Structural Inspection interleave sensing with movement and exploitation. The first four have a fixed hidden state and an observe-then-commit structure. The last two satisfy the factored-observation assumptions of Proposition~\ref{prop:factored}. Complete specifications and protocols appear in Appendices~\ref{app:envs} and~\ref{app:experimental_protocol}.

Tiger \citep{kaelbling1998} has two hidden states, one listen action of accuracy $0.85$ and cost $1$, and terminal door choices paying $+10$ or $-100$. Diagnosis has four possible conditions, two binary tests of accuracy $0.80$ and cost $1$, and a terminal diagnosis paying $+10$ or $-50$. Each test partitions the conditions differently. Bandit has four arms and a single hidden variable identifying the best arm. An inspection targets one arm, has accuracy $0.80$, and costs $0.5$. A terminal pull pays $+10$ for the best arm and $+1$ otherwise. Tileworld hides one target on a $6{\times}6$ grid. Its six noisy binary scans query bits of the target's row or column index, each with accuracy $0.80$ and cost $1$. Declaring a cell pays $+10$ or $-50$.

Our RockSample variant \citep{smith2004} has known rock positions and hidden binary qualities. Checking is noisy, movement changes sensor range, and sampling and exit yield reward. It differs materially from the original benchmark. Every action costs $0.5$, maps are our own, returns are undiscounted with an episode cap, exit is available from any cell, hidden rock qualities remain fixed after sampling, and sensor accuracy starts at $0.95$ and reaches chance at finite range. Sample history is observable and prevents a collected rock from yielding a second reward. These choices support the fixed-hidden-state analysis and make sensing use directly interpretable, but absolute returns are not comparable to published RockSample values. Appendix~\ref{app:experimental_protocol} gives the six differences in full. We evaluate RS[5,3], RS[7,4], RS[7,8], and RS[11,11], where the pair denotes grid width and rock count. Hidden-state counts exclude observable position.

Structural Inspection places $N$ components with fixed hidden nominal/faulty states on a grid. The agent moves, chooses a visual test (accuracy $0.70$, cost $0.5$) or detailed test (accuracy $0.90$, cost $2$), and makes an irreversible declaration for each component. Missing a fault costs $50$, while correct declarations pay $2$ or $5$. We evaluate $N{=}8$ and $N{=}16$, the latter with $65{,}536$ hidden configurations. This is a synthetic benchmark with a known model and a shared hand-coded search-leaf rule, not a deployed inspection study. Two additional boundary cases, the mild-penalty two-state Testbed and Navigation without separate observation actions, are reported in Appendices~\ref{app:testbed} and~\ref{app:navigation}.

Unless a battery states otherwise, the canonical seeds are $\{42,123,456,789,1024\}$. The main EFE comparison uses $1{,}000$ episodes per seed on Tiger, Diagnosis, and Bandit, and $500$ on Tileworld and Inspection. RockSample protocols appear with their tables. Means aggregate episodes across seeds, and main-table uncertainty is the SE of per-seed means. Pairwise comparisons use seed-level Welch $t$-tests with Holm--Bonferroni correction for the stated comparison families. Pooled episode-level statistics and bootstrap intervals are supplementary (Appendix~\ref{app:stats}). A failure to reject a difference is not called equivalence. Equivalence is tested separately with fixed-margin TOST below.

The budget battery generally uses $100$ episodes per seed and evaluation cell, where a cell is one weight at one seed. Tileworld uses $40$ episodes per usage-curve cell and $30$ per scale-collapse cell, and Inspection-$N{=}8$ uses $30$ per usage-curve cell. Different batteries have different episode counts and seeding rules, so their estimates are not interchangeable. The frontier study below explicitly addresses a consequential cross-battery difference. Appendix~\ref{app:repro} records the full reproducibility checklist, including all exceptions.

\section{Experiments on Shadow Prices and Sensing Budgets}
\label{sec:budget_experiments}

We first test the two operational properties of the price calibration, scale transfer and online re-adaptation. We then compare reward with unconstrained and usage-constrained references, before testing calibration-derived targets on held-out seeds. Supplementary experiments examine onset thresholds, cost-denominated budgets, interleaved settings, staircase stability, and reward-irrelevant sensing. Their principal qualifications are summarized at the end of this section.

\subsection{Curve Collapse across Reward Scales}
\label{sec:collapse}

Proposition~\ref{prop:pi1} predicts that rescaling rewards and sensing costs by $\alpha$ preserves a count-usage curve when the information weight is also multiplied by $\alpha$. We evaluate reward scales $\alpha\in\{0.1,1,10\}$ on corresponding scaled weight grids. Figure~\ref{fig:collapse} shows the result in raw and normalized coordinates. Every matched $w/\alpha$ point agrees bit-exactly on Diagnosis, Bandit, Tiger, and Tileworld. Shared per-episode random streams make exact empirical agreement possible when actions coincide. The CSVs regenerated under the relative comparison tolerance of Section~\ref{sec:pi1} are byte-identical to those generated under the earlier absolute tolerance, which did not bind at these points.

The normalized grid bracket is $(0.141,0.323]$ for Diagnosis at $B{=}8$, $(3.84,8.76]$ for Bandit at $B{=}8$, and $(5.80,20.0]$ for Tileworld at $B{=}15$, unchanged across scales. Tiger is a necessary negative control. Its usage is constant at $4.37$ across the tested normalized range $[0,20]$, so collapse is trivial and $B{=}4$ is unbracketed. Equivariance preserves a curve but does not ensure that it offers a useful sensing dial. Table~\ref{tab:collapse-breadth} and Appendix~\ref{app:scale_details} give the complete breadth check.

% FIGURE_COLLAPSE

\subsection{Online Re-adaptation after a Reward Rescale}
\label{sec:dual_multiseed}

We compare the projected decaying-step controller with its reset-on-shift extension on Diagnosis at target $B{=}8$. Both use $\eta_0{=}0.05$, decay $\delta{=}0.02$, and a tenfold reward-and-cost rescale at episode $200$ of a $400$-episode run. Ten shared controller seeds support a paired comparison. The extension watches a $20$-episode usage window. Once the step has decayed below half its initial value, a deviation exceeding both three window standard errors and $0.5$ observations restarts its decay clock. Appendix~\ref{app:dual_details} specifies the detector and steady-state summaries.

Recovery occurs when a rolling mean computed only from post-rescale episodes stays within $\pm1$ of $B$ for $20$ consecutive episodes. The comparison uses the restricted time $\min(T,180)$ for every seed, where $180$ is the last post-rescale index at which that hold can be checked. A seed that never completes it receives $180$. Decay-only averages $157.3$ episodes and reset-on-shift $51.2$, a descriptive ratio of $3.07$. The paired difference is $106.1$ episodes with a Student-$t$ $95\%$ interval $[96.8,115.4]$. All ten paired differences favor resetting. Reset-on-shift recovers on $10/10$ seeds and decay-only on $9/10$, so the latter's unrestricted recovery distribution is right-censored. The conditional mean among its nine recovered seeds is not the comparison estimand.

These measurements come from the corrected per-episode-seeded runner. The earlier runner left the rescaled environment unseeded and reported a ratio of $2.65$ with all decay-only seeds recovering. The correction preserves the qualitative advantage while revealing slower and incomplete recovery for decay-only. Figure~\ref{fig:dualmultiseed} plots the trajectories. The mid-run rescale is nonstationary, and the reset variant restarts its step schedule, so neither experiment is covered by Proposition~\ref{prop:pi5}'s stationary convergence guarantee. This is empirical re-adaptation evidence on one shift scenario.

% FIGURE_DUAL

\subsection{Unconstrained Reward Reference}
\label{sec:sarsop}

To determine whether $w{=}1$ is reward-competitive at its natural operating point, we compare EFE with SARSOP \citep{kurniawati2008} on Tiger, Diagnosis, and Bandit. The original APPL solver is built from source, each environment is exported through a common model path, and the solve uses precision $10^{-3}$ at $\gamma{=}0.999$. The returned alpha-vector policy is evaluated in the same simulator as EFE at five seeds and $500$ episodes per seed (Table~\ref{tab:sarsop}). SARSOP is a near-optimal reference, not a certified exact solution to the reported Monte Carlo comparison. Build and export details appear in Appendix~\ref{app:sarsop_build}.

% TABLE_SARSOP

Tiger's policies take identical actions on every evaluated episode. Diagnosis and Bandit show small reward differences, but non-significance alone would not establish equivalence. We therefore use two one-sided tests (TOST) \citep{schuirmann1987} with a reward margin equal to one observation's cost, $\pm1.0$ on Tiger and Diagnosis and $\pm0.5$ on Bandit. These operational margins were chosen after an earlier canonical-seed comparison had been inspected, so the five-seed TOST is retrospective. The additional fifteen seeds in the $n{=}20$ robustness study were fixed after the margin and before that study ran.

The primary TOST uses unpaired Welch tests on per-seed means. All three environments reject both one-sided nulls, including after Holm correction across environments, with $p_{\mathrm{TOST}}=0.0010$, $0.0458$, and $0.0130$ for Tiger, Diagnosis, and Bandit. Their $90\%$ intervals lie inside the specified margins. The $n{=}20$ study retains equivalence, with respective $p$-values $4.5{\times}10^{-10}$, $6.8{\times}10^{-7}$, and $6.0{\times}10^{-11}$. The nominal seed lists are shared, but differing policies consume random draws differently. A paired five-seed sensitivity analysis confirms that covariance is positive here and the unpaired analysis is conservative. Tiger's paired differences are identically zero. Diagnosis and Bandit give paired $p_{\mathrm{TOST}}$ of $0.016$ and $3.6{\times}10^{-5}$. Appendix~\ref{app:sarsop_details} gives intervals, mechanics, and archive filenames. The conclusion is equivalence within these fixed margins on these three environments, rather than universal near-optimality of $w{=}1$.

\subsection{A Usage-Constrained Reference}
\label{sec:cpomdp}

A literal usage penalty supplies a different reference from an information bonus. We replace observation reward $-\mathrm{cost}_j$ by $-(\mathrm{cost}_j+\lambda)$ and solve with SARSOP over a $12$-point penalty grid in $[0,100]$. This estimates policies maximizing $\mathbb E[R]-\lambda\mathbb E[U]$ over the POMDP policy space, with $\gamma{=}0.999$ and five seeds of $300$ episodes per penalty. The penalty is the Lagrangian price of the usage cap in Equation~\ref{eq:budgeted}. The information weight $w$ is not that multiplier.

At budget $B$, the reference is the feasible envelope of the sampled policy points, obtained by maximizing their mixture reward subject to mixture usage at most $B$. A solution uses at most two sampled policies. Policies below the target remain eligible, and the reference is flat at the largest sampled reward once the cap is slack. This is a linear program over measured points, not equal-usage interpolation. The solver is approximate, the penalty grid is finite, and reward and usage are Monte Carlo estimates. Consequently the envelope is neither a certified population bound nor the full constrained frontier. We do not assume that finite-state strong-duality results \citep{altman1999} automatically extend to the continuous belief MDP. Appendix~\ref{app:cpomdp_details} discusses these limitations and the diagnostic interpolation retained in the CSV.

% TABLE_CPOMDP

Table~\ref{tab:cpomdp} evaluates the reference at EFE's endogenous usage $B_{\mathrm{EFE}}$. The highest-reward sampled reference policy is already feasible on each environment, so the cap does not bind. Its estimated gap from EFE is zero on Tiger, $-0.041$ on Diagnosis, and $-0.005$ on Bandit, small relative to the estimates' standard errors. Negative estimated gaps can arise from sampling or approximation error and do not establish reward above the population optimum. The usage-matched Planning+IG grid weights actually run, $0.00$, $0.45$, and $1.07$, give rewards bit-identical to EFE here. This is an empirical plateau result rather than an extension of the equivalence theorem from $w{=}1$ to arbitrary weights. The reference remains limited to these small environments, with no corresponding optimality-gap estimate for the larger RockSample and Inspection instances.

\subsection{Calibration-Derived Targets on Held-Out Seeds}
\label{sec:budget_frontier}

The preceding comparison uses whatever sensing the weight-one policy happens to buy. We now test targets independent of that weight-one usage, chosen by a predeclared rule from each calibration curve. These targets stand in for application requirements but are not themselves application-specified. The study separates whether the procedure meets expected usage from whether doing so maximizes reward under a cap. Its original design was recorded in the producer's docstring before results were inspected (\texttt{experiments/\allowbreak run\_\allowbreak budget\_\allowbreak frontier.py}).

Each core environment is calibrated on the canonical five seeds, $100$ episodes per seed, using the $16$-point weight grid and horizons of Appendix~\ref{sec:stairs}. We select budgets at $0.25$, $0.5$, and $0.75$ of the measured range $[U_{\min},U_{\max}]$, at the midpoint of its largest jump, and at half of $U_{\min}$. The last is deliberately unattainable and must be declared so. Definition~\ref{def:pi3} supplies the crossing bracket and endpoint-mixture probability $q$. All selection uses calibration data alone, and evaluation uses held-out seeds $\{7,8,9,10,11\}$ with $100$ episodes each. Tiger's gap budget duplicates its middle-range budget. Bandit's curve is nonmonotone, so the stipulated last-crossing convention determines its bracket.

The original within-family comparator is the highest-reward single grid member whose calibration usage is at most $B$. Two additional LP comparators were introduced after inspecting the first results, and that amendment is disclosed in the producer. One maximizes calibration reward over mixtures with usage equal to $B$, and one does so with usage at most $B$. Each uses at most two calibration-grid weights and is evaluated on held-out seeds. Their scope is the sampled grid. Selecting calibration usage below a cap does not guarantee that the population or held-out usage meets it.

Table~\ref{tab:budget_frontier_main} gives the principal held-out results. All three below-range targets are declared unattainable. At all twelve attainable targets, the selected endpoint mixture meets usage to within $0.13$ observations. Errors are within one held-out SE except Bandit's gap target, where an error of $0.09$ is about $1.3$ SE. No single grid weight approximates these targets as closely. Endpoints miss by $0.3$--$1.1$ observations on Tiger and $0.2$--$3.6$ elsewhere. The selected feasible single member leaves $0.2$--$4.7$ observations unspent on held-out data. Its usage is below the target in all twelve rows, an observed result in addition to the calibration selection fact. Conversely, the endpoint mixture's held-out mean exceeds $B$ on three rows, by at most $0.13$. An expected-use target provides neither a hard per-episode cap nor finite-sample cap feasibility.

% TABLE_FRONTIER_COMPACT

Meeting a target can buy useful information or impose an avoidable reward cost. On Diagnosis at $B{=}7.72$, $9.53$, and its gap target $7.84$, the mixture earns $-2.14$, $-2.11$, and $-2.06$ against the feasible member's $-2.86$. On Tiger, reward-only planning uses $4.32$ observations and earns $5.68$ on this held-out stream. Spending more to meet the three increasing targets reduces reward to $5.36$, $4.96$, and $4.65$. Bandit likewise loses reward at its two larger interior targets. The endpoint procedure also need not select the best mixture meeting a target. At Bandit's gap budget, the calibration LP mixes $w{=}0$ and $w{=}0.72$ and earns $5.94\pm0.19$ on held-out seeds, against the endpoint mixture's $5.62\pm0.21$. Its usage error is larger, $0.18$ rather than $0.09$. The chosen endpoint mixture therefore cannot stand for every mixture in the family, and its deficit at this budget cannot be attributed to planning depth alone.

A valid constrained-reward comparison must also control the evaluation stream. The original reference points were evaluated on canonical seeds with once-per-seed seeding, whereas the family uses held-out per-episode seeds. Their direct reward comparisons, retained in Appendix~\ref{app:frontier_details}, are descriptive cross-stream comparisons. The difference matters on Tiger. Its unpenalized reference earns $5.02$ on the canonical stream but $5.68$ on the held-out stream, removing the mixture's apparent excess over the reference at the smallest budget.

We therefore added a same-stream comparison after inspecting the original results. Every saved reference policy was first checked against its committed canonical estimate to $10^{-9}$, then re-evaluated on the held-out seeds with the family's per-episode seeding, and the feasible envelope was recomputed. The final column of Table~\ref{tab:budget_frontier_main} pairs each seed's mixture reward with its reference reward on the same episode seeds. Reference support and mixture weights are fitted on those same held-out seed means and held fixed for the SE. The resulting Student-$t$ intervals are nominal pointwise plug-in summaries. They omit support reselection, uncertainty in mixture weights and cap feasibility, and multiplicity across budgets. They do not establish coverage for the population constrained frontier.

Under that fitted comparison, negative $95\%$ intervals exclude zero at nine of twelve attainable budgets, or eight of eleven distinct budgets after the Tiger duplicate is removed. The largest measured shortfall is $1.32$ on Bandit, about one fifth of its reference reward. Three intervals include zero, at Diagnosis's smallest and gap targets and Bandit's smallest target. Re-evaluation removes the stream mismatch but leaves selection bias from maximizing over noisy reference points.

Tiger requires an additional rare-event qualification. Neither the family nor any policy in the fitted reference support makes a wrong commit in the $500$ held-out episodes. Its paired gaps therefore measure only extra listening. A wrong commit costs $110$ relative to a correct one, and the canonical reference error rate is $0.6\%$, equivalent to about $0.66$ expected reward. If the reference retained that rate while the family never erred, the smallest-budget shortfall would reverse and the middle-budget shortfall would nearly vanish. Only the largest Tiger shortfall is robust to that sensitivity calculation. The nominal intervals do not capture this rare-event uncertainty.

This study supports expected-usage calibration within the tested family, including gap targets, while quantifying its limits as a reward optimizer. If the engineering requirement is to spend a specified expected sensing amount for coverage or verification that reward does not represent, the endpoint mixture directly addresses it. If the requirement is only an upper limit and reward is the objective, the feasible mixture or constrained reference is the relevant comparator. The complete three-panel table, selection details, historical reference correction, and paired archives are retained in Appendix~\ref{app:frontier_details}.

\subsection{Breadth and Limits of the Budget Instrument}
\label{sec:budget_breadth_summary}

The supplementary studies test distinct parts of the calibration story. Positive-threshold two-state instances place the observed onset within one refined grid step, $3\%$, above the closed-form threshold, while standard negative-threshold instances observe already at $w{=}0$ (Appendix~\ref{sec:prop2exp}). Heterogeneous test costs show that observation-count and sensing-cost targets are not interchangeable when the test mix changes with $w$ (Appendix~\ref{sec:cost_budget}). RockSample and Inspection exercise calibration in interleaved settings, subject to the fixed-hidden-state scope and search limitations already stated (Appendix~\ref{sec:interleaved_budget}). The measured staircases and atlas record operating points rather than a universal predictive formula for the price (Appendices~\ref{sec:stairs} and~\ref{app:w_atlas}). Bootstrap bracket stability measures resampling stability of the grid selection, not a confidence interval for the population price.

A budget also cannot ensure that sensing concerns reward-relevant uncertainty. In Distractor Diagnosis, ordinary Planning+IG begins buying an independent, reward-irrelevant test between $w{=}3.5$ and $4.0$, and roughly a third of usage goes to it at large weights. The weight-one agent avoids it on this instance only because it lies below onset. Scoring information gain on reward-equivalence classes removes distractor tests at every swept weight, without changing belief propagation, but can still over-buy task-relevant information. The reward-equivalence partition is read from the supplied commit reward matrix. The IDS adaptation also buys distractor tests because its zero-denominator fallback uses raw state entropy. This is a limitation of that implementation, not a general impossibility for reward-aware sensing. Appendix~\ref{sec:distractor} retains the paired curves, corrected tests, implementation explanation, and reward costs. Together these results distinguish how much to sense from what to sense about.

\section{Results}
\label{sec:results}

The second battery locates the canonical EFE weight within the policy family and compares it with other planners. It provides empirical context for the equivalence result, whose validity itself rests on the stated assumptions and proof. We retain the main comparisons here and place detailed per-environment mechanisms, robustness studies, and additional tables in Appendix~\ref{app:extended_efe_results}.

\subsection{Core Environments}

Table~\ref{tab:main} contrasts $w{=}1$ EFE with myopic reward maximization, same-horizon reward-only Planning, and Planning+IG tuned for success within its own class. Reward-maximizing weights are considered separately below. On Tiger, EFE and Planning produce identical trajectories. With one observation action, the information term never changes the reward-based action choice at these stakes. This is a measured property of this instance rather than a general consequence of having a single sensor.

% TABLE_CORE

On Diagnosis, EFE raises success from $88.8\%$ to $97.2\%$ relative to Planning and improves reward from $-2.57$ to $-1.37$. Seed-level Welch tests give $p=4.0{\times}10^{-8}$ for success and $p=0.002$ for reward. On Bandit it raises success from $71.0\%$ to $86.9\%$ and reward from $5.75$ to $6.27$, with $p=2.1{\times}10^{-5}$ and $p=0.001$. These differences survive the stated Holm correction. Thus EFE Pareto-dominates same-horizon reward-only Planning on these two environments. Success-tuned Planning+IG reaches higher accuracy, $99.2\%$ and $98.9\%$, but pays for more sensing and attains lower reward. EFE supplies an untuned trade-off, not the best value of every metric.

Two baselines limit any attribution of these outcomes solely to explicit information gain. Posterior-vote, which draws states from the belief and votes over their optimal actions without an information term, is not significantly different from EFE on any Bandit metric. Its reward is nominally higher, $6.32$ against $6.27$ ($p=0.58$). It also beats EFE on Testbed, but performs worse on Diagnosis and Tiger. The Epistemic-only ablation commits immediately at chance-level success on all three core environments and Tileworld. Removing the pragmatic term therefore causes degenerate non-exploration in this implementation. Full baselines, pooled effect sizes, and seed-level tests appear in Appendices~\ref{app:full_tables}, \ref{app:stats}, and~\ref{app:core_details}. Reward effects against Planning are small at the episode level on Diagnosis and Bandit despite statistically detectable mean differences. We do not interpret the mechanically inflated effect sizes computed from only five seed means.

POMCP provides a search-based comparison, reported in Appendix~\ref{app:pomcp}. Its simulation count per decision is a computational budget distinct from the sensing budget $B$. The main POMCP battery is simulation-matched, and the Tiger wall-clock check is a separate matched-compute result. The original POMCP comparisons use untuned exploration constants. A separate sweep selects constants on tuning seeds under a recorded rule, but the canonical five-seed sweep had already been inspected before that rule was written. The Tiger tuning result is therefore retrospective. Neither the POMCP comparison nor the component ablation isolates directed observation selection as the cause of every performance difference. These qualifications remain necessary when comparing search architectures, even when the information-weighted family performs well.

\subsection{Where the Information-Unit Weight Sits}
\label{sec:pareto}

We sweep $w\in\{0.01,0.1,0.5,1,2,5,10,20,50,100,200\}$ on Tiger, Diagnosis, Bandit, Testbed, and Tileworld. Figure~\ref{fig:pareto} shows reward against success. The $w{=}1$ Planning+IG and EFE markers coincide under Proposition~\ref{prop:equivalence} and shared tie-breaking. The empirical question is where that common policy lies relative to other weights, not whether equivalent implementations differ.

% FIGURE_PARETO

The weight-one policy is within the reward-maximizing run of sampled weights on Tiger ($0.01$ through $20$), Diagnosis ($0.5$ through $20$), and Bandit ($1$ and $2$). These are sampled plateaus, not claims about every untested weight between grid points. Testbed instead maximizes reward across sampled weights from $0.01$ through $0.5$. Its weight-one policy buys $6.4$ percentage points of success at a reward cost of $0.11$. On Tileworld, $w{=}20$ Pareto-dominates $w{=}1$, with reward $-18.95$ against $-21.53$ and success $91.8\%$ against $72.0\%$. Thus the canonical weight attains the sampled reward maximum on three of five environments and is suboptimal on two.

Success favors larger weights throughout this sweep, with its maximum at the upper grid endpoint $200$ on all five environments. The reward cost need not grow smoothly. On Tiger and Diagnosis, the tested weights from $1$ up to the point below $50$ leave reward unchanged, and the step to $50$ incurs the measured cost. The grid does not exclude a cheaper untested intermediate policy. The advantage of $w{=}1$ here is that it selects a useful default without environment-specific search under the stated reward convention. It does not remove dependence on units, promise best accuracy, or replace calibration when a particular sensing target is required. Appendix~\ref{app:pareto_details} retains the detailed trade-offs, and Appendix~\ref{app:nat_canonical} compares the implementation's bit-unit default with the exactly nat-canonical weight.

\subsection{Tileworld}
\label{sec:tileworld}

Tileworld shows why horizon and observation structure must accompany a scaling claim. At $6{\times}6$ and $H{=}2$, EFE and Planning are not significantly different in reward ($-21.53$ against $-21.39$, $p=0.82$) or success ($72.0\%$ against $73.8\%$, $p=0.17$). EFE uses fewer scans, $14.75$ against $15.64$ ($p=8.8{\times}10^{-5}$). These are failures to reject reward and success differences, not TOST equivalence. The success-tuned $w{=}50$ agent reaches $97.3\%$ accuracy but uses $32.25$ scans and earns lower reward. Table~\ref{tab:tileworld} and Figure~\ref{fig:tw_comparison} retain the full comparison and an illustrative trajectory in Appendix~\ref{app:tileworld_details}.

In a separate fixed-$H{=}2$ sweep, Planning takes no scans at $8{\times}8$ and achieves $1.4\%\pm0.3$ percentage points success, while EFE achieves $69.4\%\pm1.0$. The short horizon cannot value enough scans to repay their costs before committing on a large grid. The contrast is therefore with a non-scanning finite-horizon agent, not with reward-only planning at arbitrary depth. At $4{\times}4$ and $6{\times}6$ the two are within one SE, with no formal test for that sweep. The advantage appears at a horizon-and-scale threshold rather than widening uniformly with size. Its weighted comparators use fixed $w{=}100$, not class-specific tuning. Random and overlapping scan partitions sharply reduce all agents' absolute performance and do not separate EFE from Planning on reward or success, although EFE uses fewer scans in all three modes. Appendix~\ref{app:tileworld_details} preserves these robustness results and their scope.

\subsection{RockSample}
\label{sec:rocksample}

RockSample tests the information weight within a shared factored belief-space search and a shared greedy leaf rule. On RS[5,3], EFE ($16.47$) and the $w{=}5$ comparator ($16.41$) are not significantly different. On RS[7,4], EFE earns $15.96$ against $14.56$ and leads the weighted family ($p<10^{-5}$). On RS[7,8], EFE ($21.85$) and $w{=}5$ ($23.76$) are not significantly different after correction, and both outperform reward-only Planning ($12.31$). Raising the weight from $5$ to $10$ significantly hurts on two of these instances and does not significantly help on the third. Only weights $1$, $5$, and $10$ were compared here, so this is not a dense reward-tuning claim.

The important counterexample is RS[11,11]. EFE and Planning both earn $13.23$, while a standalone approach-then-check heuristic earns $28.66\pm0.83$. That heuristic also beats both weighted agents on RS[7,8]. Larger state space therefore does not itself increase the value of information weighting. At feasible search depths, distant useful checks lie beyond the tree, and the shared leaf rule values untouched rocks at zero under the uniform prior. It also charges a phantom trip to the east edge before exit, although this variant allows exit anywhere, biasing depth-capped continuations toward early exit. The cost is shared across the weighted family but limits absolute performance. Deeper POMCP does not repair this instance in the tested configurations. Behavioral traces show many distant checks, and reward declines from $12.94$ at horizon $10$ to $-54.78$ at horizon $40$. Appendix~\ref{app:rocksample_main_details} retains the main table, leaf-rule disclosure, depth check, and heuristic comparison, alongside the extended solver analysis in Appendix~\ref{app:rocksample}.

\subsection{Structural Inspection}
\label{sec:inspection}

Inspection makes the reward-accuracy trade-off explicit. At $N{=}8$, EFE improves accuracy from Planning's $73.0\%$ to $91.1\%$ while reward falls from $-17.85$ to $-20.95$. At $N{=}16$, it improves accuracy from $78.2\%$ to $90.8\%$, with reward $-45.80$ against $-45.42$, a non-significant difference ($p=0.57$). The $w{=}5$ agent attains higher accuracy, $97.3\%$, but uses $38.3$ tests against EFE's $28.1$ and earns $-54.74$. These results demonstrate feasible factored search on a synthetic known-model problem with a hand-coded leaf rule. They do not validate industrial deployment or show that EFE uniquely optimizes either reward or accuracy. Table~\ref{tab:inspection} and the qualified two-state interpretation appear in Appendix~\ref{app:inspection_details}.

Across the battery, the canonical weight is useful when the reward-only policy leaves valuable, representable sensing opportunities unused. Its benefit is bounded by the search horizon, sensor structure, reward convention, and objective of the comparison. It matches the sampled reward maximum on three of five swept environments, loses to other weights on two, and gains nothing over Planning on the largest RockSample instance. These limits matter alongside its positive results on Diagnosis, Bandit, and the largest fixed-horizon Tileworld grid.
